\chapter{Complete \(\mu_1\)--\(\mu_5\) Pipeline Code}

This appendix contains annotated core implementation code.

\section{Stage 1: Normalization}

\begin{minted}[fontsize=\footnotesize]{rust}
use oxigraph::store::Store;
use oxigraph::model::Graph;

pub fn normalize(input: &str) -> Result<Graph, NormalizeError> {
    // Parse Turtle
    let store = Store::new();
    let graph = parse_turtle(input, &store)?;

    // Validate SHACL
    let report = validate_shacl(&graph)?;
    if !report.is_conformant() {
        return Err(NormalizeError::ShaclViolation(report));
    }

    // Check entropy
    let entropy = calculate_entropy(&graph);
    if entropy > 20.0 {
        warn!("High entropy: {} bits (possible incompleteness)", entropy);
    }

    // Verify closure
    verify_closure(&graph)?;

    Ok(graph)
}

fn verify_closure(graph: &Graph) -> Result<(), NormalizeError> {
    for reference in get_references(graph) {
        let definition = query_definition(graph, &reference)?;
        if definition.is_none() {
            return Err(NormalizeError::UndefinedReference(reference));
        }
    }
    Ok(())
}
\end{minted}

\section{Stage 2: Extraction}

\begin{minted}[fontsize=\footnotesize]{rust}
use oxigraph::sparql::{Query, QueryResults};
use oxigraph::model::Graph;

pub fn extract(graph: &Graph) -> Result<ExtractedData, ExtractError> {
    let mut data = ExtractedData::new();

    // Run all extraction patterns
    for pattern in STANDARD_PATTERNS {
        let results = execute_query(graph, &pattern.query)?;
        data.merge(results);
    }

    Ok(data)
}

fn execute_query(graph: &Graph, query_str: &str) -> Result<QueryOutput, ExtractError> {
    let query = Query::parse(query_str, None)?;
    match query.evaluate(graph) {
        Ok(QueryResults::Graph(g)) => Ok(QueryOutput::Graph(g)),
        Ok(QueryResults::Solutions(s)) => Ok(QueryOutput::Solutions(s)),
        Err(e) => Err(ExtractError::QueryFailed(e)),
        _ => Err(ExtractError::UnexpectedResult),
    }
}

const STANDARD_PATTERNS: &[Pattern] = &[
    Pattern {
        name: "Entity Discovery",
        query: r#"
            PREFIX rdf: <http://www.w3.org/1999/02/22-rdf-syntax-ns#>
            CONSTRUCT { ?entity a ?type }
            WHERE { ?entity a ?type }
        "#,
    },
    // ... more patterns
];
\end{minted}

\section{Stage 3: Emission}

\begin{minted}[fontsize=\footnotesize]{rust}
use tera::{Tera, Context};

pub fn emit(data: &ExtractedData, rules: &[Rule]) -> Result<Vec<Artifact>, EmitError> {
    let mut artifacts = Vec::new();
    let tera = compile_templates()?;

    for rule in rules {
        let context = prepare_context(data, rule)?;
        let output = tera.render(&rule.template, &context)?;
        artifacts.push(Artifact {
            path: rule.output.clone(),
            content: output,
        });
    }

    Ok(artifacts)
}

fn prepare_context(data: &ExtractedData, rule: &Rule) -> Result<Context, EmitError> {
    let mut context = Context::new();
    context.insert("entities", &data.entities);
    context.insert("relationships", &data.relationships);
    context.insert("config", &rule.config);
    Ok(context)
}

fn compile_templates() -> Result<Tera, EmitError> {
    let mut tera = Tera::default();
    tera.add_raw_templates(vec![
        ("rust_struct.rs", include_str!("templates/rust_struct.rs")),
        ("typescript_interface.ts", include_str!("templates/typescript_interface.ts")),
    ])?;
    Ok(tera)
}
\end{minted}

\section{Stage 4: Canonicalization}

\begin{minted}[fontsize=\footnotesize]{rust}
use std::process::Command;

pub fn canonicalize(artifact: &Artifact) -> Result<String, CanonicalizeError> {
    let extension = artifact.extension();
    let formatted = match extension {
        "rs" => canonicalize_rust(artifact)?,
        "ts" | "tsx" => canonicalize_typescript(artifact)?,
        "py" => canonicalize_python(artifact)?,
        "yaml" | "yml" => canonicalize_yaml(artifact)?,
        _ => artifact.content.clone(),
    };
    Ok(formatted)
}

fn canonicalize_rust(artifact: &Artifact) -> Result<String, CanonicalizeError> {
    let output = Command::new("rustfmt")
        .arg("--emit")
        .arg("stdout")
        .stdin(std::process::Stdio::piped())
        .stdout(std::process::Stdio::piped())
        .spawn()?
        .stdin
        .unwrap()
        .write_all(artifact.content.as_bytes())?;

    // ... read output
    Ok(formatted)
}
\end{minted}

\section{Stage 5: Receipt Generation}

\begin{minted}[fontsize=\footnotesize]{rust}
use ed25519_dalek::{Keypair, Signature, Signer};
use blake3::Hash;

#[derive(Debug, Clone)]
pub struct Receipt {
    pub artifact_hash: Hash,
    pub spec_hash: Hash,
    pub signature: Signature,
    pub timestamp: chrono::DateTime<chrono::Utc>,
    pub metadata: Metadata,
}

pub fn generate_receipt(
    artifact: &str,
    spec_hash: Hash,
    keypair: &Keypair,
) -> Receipt {
    let artifact_hash = blake3::hash(artifact.as_bytes());
    let combined = combine_hashes(spec_hash, artifact_hash);
    let signature = keypair.sign(&combined);

    Receipt {
        artifact_hash,
        spec_hash,
        signature,
        timestamp: chrono::Utc::now(),
        metadata: Metadata {
            version: env!("CARGO_PKG_VERSION").to_string(),
            generator: "ggen".to_string(),
        },
    }
}

fn combine_hashes(h1: Hash, h2: Hash) -> [u8; 64] {
    let mut combined = [0u8; 64];
    combined[..32].copy_from_slice(&h1.as_bytes());
    combined[32..].copy_from_slice(&h2.as_bytes());
    combined
}

pub fn verify_receipt(receipt: &Receipt, public_key: &PublicKey) -> bool {
    let combined = combine_hashes(receipt.spec_hash, receipt.artifact_hash);
    public_key.verify(&combined, &receipt.signature).is_ok()
}
\end{minted}
